\documentclass[12pt, a4paper]{article}
\usepackage[utf8]{inputenc}
\usepackage{amsmath}
\usepackage{amsfonts}
\usepackage{geometry}
\usepackage{fancyhdr}
\usepackage{verbatim}
\usepackage{booktabs}

% Page geometry to maximize space for the problem statement
\geometry{left=1in, right=1in, top=1in, bottom=1in}

\begin{document}

% Problem Header
\begin{center}
    {\huge \textbf{Amicable Numbers}} \\
    \vspace{0.2cm}
    \textbf{Time Limit:} 2.0s \quad | \quad \textbf{Memory Limit:} 512 MB
\end{center}

\vspace{0.5cm}

% Problem Description
\section*{Description}
Some numbers form bonds with others, not just any bond, but a very specific kind. They don't shake hands, but they share something quietly through their parts.

\subsection*{Clue: The Breakdown}
Take a number $n$. Break it down. Not all the way, just the pieces that evenly go into it (excluding the number itself). These are called the \textbf{proper divisors}. Now, sum those parts. What you get might lead you to someone familiar.

\subsection*{Clue 1: A Special Connection}
Try it with $220$:
\begin{itemize}
    \item Proper divisors: $1, 2, 4, 5, 10, 11, 20, 22, 44, 55, 110$
    \item Sum: $284$
\end{itemize}

Now try the same with $284$:
\begin{itemize}
    \item Proper divisors: $1, 2, 4, 71, 142$
    \item Sum: $220$
\end{itemize}

That is a friendship. Two numbers, mutually pointing to each other. They are \textbf{amicable}.

\subsection*{Clue 2: Not All Numbers Are Lucky}
Try $100$:
\begin{itemize}
    \item Proper divisors: $1, 2, 4, 5, 10, 20, 25, 50 \rightarrow$ Sum is $117$
\end{itemize}
Now for $117$:
\begin{itemize}
    \item Divisors are $1, 3, 9, 13, 39 \rightarrow$ Sum is $65$
\end{itemize}
No cycle, no connection. Just numbers doing their own thing.

\vspace{0.3cm}
\noindent \textbf{Definition:}
Amicable numbers must mutually recognize each other through the sum of their proper divisors. However, the numbers should not be equal.
\begin{itemize}
    \item For example, $6$ is a perfect number (sum of divisors is $6$), so it has no friend except itself. In this problem, we do not consider $n$ amicable if its friend is $n$.
    \item $1$ does not have any proper divisors, so it has no friends either.
\end{itemize}

\section*{Task}
You will be given multiple numbers. For each one, determine if it belongs to an amicable pair.

\section*{Input}
The first line contains an integer $t$, the number of test cases. \\
The next $t$ lines each contain a single positive integer $n$.

\subsection*{Constraints}
\begin{itemize}
    \item $1 \le t \le 10^{5}$
    \item $1 \le n \le 10^{7}$
\end{itemize}

\section*{Output}
For each number $n$, output a single line:
\begin{itemize}
    \item If $n$ is part of an amicable pair, print its friend.
    \item Otherwise, print \texttt{$-1$}.
\end{itemize}
\textit{Note: The friend of $n$ is guaranteed to be less than $4\times10^{7}$.}

\section*{Sample}

\begin{table}[h]
\centering
\begin{tabular}{|p{0.45\textwidth}|p{0.45\textwidth}|}
\hline
\textbf{Input} & \textbf{Output} \\
\hline
\begin{verbatim}
3
220
284
100
\end{verbatim}
&
\begin{verbatim}
284
220
-1
\end{verbatim}
\\
\hline
\end{tabular}
\end{table}

\vspace{0.5cm}
\noindent \textit{Think of each number like a puzzle made up of smaller pieces. Some of these pieces, the proper divisors, quietly hold the key to something special. Every now and then, two numbers spot each other through these parts. When both reach out in return, that is when a true friendship is formed.}

\end{document}